\documentclass[]{uiuc_paper}
\pdfoutput=1
% single-column (default): \documentclass[]{uiuc_paper}
% two-column: \documentclass[twocolumn]{uiuc_paper}

%% =====================================================================
%%  七部分结构 —— 与分工表一一对应
%%
%%   1  Introduction        sections/01-introduction.tex        成员 1
%%   2  Literature Review   sections/02-literature-review.tex    成员 2
%%   3  Methodology         sections/03-methodology.tex          成员 3
%%   4  Results             sections/04-results.tex              成员 4
%%   5  Discussion          sections/05-discussion.tex           成员 5
%%   6  Conclusion          sections/06-conclusion.tex           成员 6
%%   7  Abstract + 统稿     sections/00-abstract.tex             成员 7
%%
%%  编译：latexmk -pdf paper.tex        （改动后重跑，让 \ref 收敛）
%%  清理：latexmk -C paper.tex
%%  说明：第 7 部分是摘要 + 全文统稿，摘要排版在标题页的摘要框里，
%%        所以正文编号是 1--6，摘要不再单列一节。
%%  页眉校徽：uiuc_paper.cls 里的 \headerlogo（默认 figs/seu-logo.png）；
%%        要换图，改那一行，或在导言区 \renewcommand{\headerlogo}{figs/xxx}
%% =====================================================================

\usepackage[toc,page,header]{appendix}
\usepackage{cleveref}
\usepackage{subcaption}
\usepackage{booktabs}
\usepackage{tabularx}
\usepackage{colortbl}
\usepackage{algorithm}
\usepackage{algpseudocode}
\usepackage{inconsolata}
\usepackage{needspace}

% ======== Standard Packages ========
\usepackage{latexsym}
\usepackage{url}
\usepackage{amssymb}
\usepackage{amsmath}
\usepackage{amsfonts}
\usepackage[utf8]{inputenc}
\usepackage{booktabs}
\usepackage{pifont}
\usepackage{multirow}
\usepackage{makecell}
\usepackage{paralist}
\usepackage{xspace}
\usepackage{adjustbox}

% ======== Common Macros ========
\newcommand{\todo}[1]{\textcolor{red}{{[$\blacksquare$ TODO: #1]}}}
\newcommand{\ie}{\textit{i.e.}\xspace}
\newcommand{\eg}{\textit{e.g.}\xspace}
\newcommand{\wo}{\textit{w/o}\xspace}
\newcommand{\etc}{\textit{etc.}\xspace}
\newcommand{\dq}[1]{``#1''}
\newcommand{\sq}[1]{`#1'}
\newcommand{\cmark}{\ding{51}}
\newcommand{\xmark}{\ding{55}}

% ======== Author Comment Macros (remove for camera-ready) ========
\newcommand{\jj}[1]{\textcolor{cyan}{\small{\bf [JJ: #1]}}}
% \newcommand{\coauthor}[1]{\textcolor{blue}{{[Coauthor: #1]}}}

% ======== Footnote without number ========
\newcommand\footnoteplain[1]{%
  \begingroup
  \renewcommand\thefootnote{}\footnote{#1}%
  \addtocounter{footnote}{-1}%
  \endgroup
}


% ======== 方法名 ========
% 全文统一用 \method 指代自己的方法，避免同一个名字出现多种写法
\newcommand{\method}{MethodName\xspace}

% ======== 配色（与 figs/ 里的图保持一致）=======
% navy = 本文方法，blues = 对比/基线，pink = 反例
\definecolor{cnavy}{HTML}{2F4A6D}
\definecolor{cbluedeep}{HTML}{506684}
\definecolor{cbluemid}{HTML}{8797B0}
\definecolor{cbluesoft}{HTML}{B7C6DC}
\definecolor{cice}{HTML}{DCE2EA}
\definecolor{cpink}{HTML}{DCC2BE}
\definecolor{cink}{HTML}{2B3F5B}
\definecolor{crose}{HTML}{A47C7A}
\colorlet{oursrow}{cbluesoft!28}   % 表格里本文方法那一行的底色

% ======== 参考文献开关 ========
% 正文里第一次出现 \citep{...} 之后，把下面这行的 false 改成 true
% （不改成 true 时不会生成 References 页，避免空页面）
\newif\ifshowbib
\showbibtrue

%% =====================================================================
%%  格式约定（成员 4/7 尤其注意，保持全文一致）
%%    · 图表标题：\caption{\textbf{一句话结论.} 设置、指标口径、平均了几次}
%%    · 表格：booktabs 三线表，\centering\small，\tabcolsep 4.8--8pt
%%    · 正文引用图表用 Figure~\ref{} / Table~\ref{}；引用文献用 \citep{}
%%    · 未完成的段落用 \todo{...} 标记，交稿前全局搜索 "TODO" 清理
%% =====================================================================

% ======== 标题 ========
% 填你们论文的题目；术语大小写由成员 7 最后统一
\title{From Advisor to Agent: Rethinking the Human--AI Relationship in Medicine}

% ======== 作者 ========
% 七位成员各一行；方括号里的数字对应下方 \affiliation 的编号
\author[1]{Member 1}
\author[1]{Member 2}
\author[1]{Member 3}
\author[1]{Member 4}
\author[1]{Member 5}
\author[1]{Member 6}
\author[1]{Member 7}

% ======== 单位 ========
\affiliation[1]{Southeast University}

% ======== 贡献标注（可选）=======
% \contribution[*]{Equal contribution}

% ======== 摘要：第 7 部分 ========
\abstract{%% =====================================================================
%%  7 / 7   Abstract  摘要 + 全文统稿                           成员 7
%%  五句话结构：背景 / 空白 / 做法 / 结果 / 结论
%% =====================================================================

Medical artificial intelligence is moving from systems that recommend to systems that act. Large language model agents now plan, call tools and complete multi-step clinical tasks, yet they are evaluated almost exclusively on accuracy: the strongest public benchmark reports a success rate below 70\%. This paper asks what changes when the unit of analysis moves from the advice to the action. We conduct a purposive review of seventeen studies spanning engineering and human-factors work, and identify four themes --- deskilling, the supervision illusion, accountability diffusion, and trust under uncertainty --- which we show to be symptoms of a single missing object: an explicit authorization boundary. Existing frameworks describe how much autonomy a system has; none specifies what a delegation must contain. We propose \emph{governed delegation}, which requires four things of any agent permitted to act clinically: a stated scope of authority, revocability while the action is in flight, auditability afterwards, and retention of unaided human skill. We then specify a $2\times2$ vignette experiment that tests the advisor-to-agent shift directly. The human--AI relationship in medicine is less a question of trust than a question of who may act, and on whose authority.
}

% ======== 元信息 ========
\date{\today}
\correspondence{\email{name@example.com}}

\begin{document}

% ======== 定理环境（按需使用）=======
\newtheorem{theorem}{Theorem}[section]
\newtheorem{proposition}[theorem]{Proposition}
\newtheorem{lemma}[theorem]{Lemma}
\newtheorem{corollary}[theorem]{Corollary}
\newtheorem{definition}[theorem]{Definition}
\newtheorem{assumption}[theorem]{Assumption}
\newtheorem{remark}[theorem]{Remark}

\maketitle

%% ---- 正文：六个部分，一人一个文件 ----
%% =====================================================================
%%  1 / 7   Introduction  引言                                   成员 1
%% =====================================================================
\section{Introduction}
\label{sec:intro}

\paragraph{Background.}
Clinical artificial intelligence (AI) has spent a decade being useful in the passive voice. Systems read images, flag physiological deterioration and rank differential diagnoses, but they return a judgement and stop. That division of labour was stable enough to be codified in regulation and in training: the clinician decided, the software advised, and the interface between them was a screen. Two developments have disturbed it. The first is capability. Large language models have been wrapped into agents that plan, call external tools, write to records and execute multi-step procedures, so that the output of a clinical system is increasingly not a recommendation but a completed action \citep{wang2025baymax}. The second is institutional. Such systems are being deployed into workflows whose volume makes case-by-case review impractical, and they are marketed on the promise that review is no longer needed.

\paragraph{Problem.}
Evaluation has not changed to match. Benchmarks measure task success, and success is reported as if it were the only property that matters. On \textsc{MedAgentBench}, the most widely used virtual electronic health record environment, the strongest model completes 69.67\% of tasks \citep{jiang2025medagentbench}. A system that fails roughly one task in three is a legitimate object of study; it is not an obvious object of delegation. Yet the literature reports the number and leaves the delegation unexamined.

Two research communities describe the same phenomenon without meeting. Engineering work asks how to build agents that act correctly, and improves architectures, memory and tool use. Human-factors work asks what happens to people who work alongside automation, and has documented complacency, skill decay and misplaced mistrust for four decades \citep{parasuraman2010}. Neither asks the question that now matters: under what conditions may an agent be authorized to act at all?

A terminological looseness compounds the problem. ``Medical agent'' is used for systems ranging from a single-turn question answerer to an autonomous ward-round simulator. Without a stable referent, claims about autonomy cannot be compared, and a gain measured at one level is read as evidence about another.

\paragraph{Contributions.}
This paper makes four contributions.
\begin{itemize}
  \item We give a scoping definition of a medical agent --- planning, tool use and multi-step execution, as distinct from single-turn question answering --- and use it to bound the corpus.
  \item We conduct a purposive review of seventeen studies across six databases, placing engineering and human-factors sources in one corpus and coding each for its locus of action and level of autonomy.
  \item We identify four themes --- deskilling, the supervision illusion, accountability diffusion, and trust under uncertainty --- and show that all four are symptoms of one missing object: an explicit authorization boundary.
  \item We propose \emph{governed delegation}, which specifies four requirements for any agent permitted to act, and set out a $2\times2$ vignette protocol that tests the advisor-to-agent shift directly.
\end{itemize}

\paragraph{Roadmap.}
Section~\ref{sec:litreview} reviews the three literatures and states the gap. Section~\ref{sec:method} describes the review protocol and the experimental design. Section~\ref{sec:results} reports the coding and the four themes. Section~\ref{sec:discussion} develops governed delegation and answers two objections, and Section~\ref{sec:conclusion} concludes.

%% =====================================================================
%%  2 / 7   Literature Review  文献综述                         成员 2
%%  三条主线 + 研究空白
%% =====================================================================
\section{Literature Review}
\label{sec:litreview}

\paragraph{Stream 1: from decision support to agentic systems.}
\citet{topol2019} framed the field as a convergence of human and artificial intelligence, in which deep learning would improve diagnosis without displacing the clinician. That framing held while clinical systems remained classifiers. The survey of \citet{wang2025baymax} marks the change: it organizes medical agents by level of involvement and by paradigm, and makes explicit that the object of study is now a system that acts over time rather than a function that returns a label. Concrete instances followed quickly. Agent Hospital simulates an entire hospital in which patient agents and doctor agents evolve together, so that policies can be improved without risk to real patients \citep{li2024agenthospital}. \citet{zhuang2025learning} search over medical agent architectures and report task success of 91.0\% on their benchmark, rising to 98.0\% with a memory module. \citet{jiang2025medagentbench} build a virtual health record in which agents must read, write and act. \citet{schmidgall2024agentclinic} place agents in simulated consultations, \citet{gao2026medexagent} extend evaluation to clinical examination, and \citet{ma2026activemedagent} show that on 175 cases an active agent reaches the correct diagnosis while requesting fewer examinations than a passive one.

The limitation is uniform across the stream: every one of these studies measures what the system can do. None measures who is permitted to let it.

\paragraph{Stream 2: autonomy, oversight and accountability.}
A second literature treats autonomy as a graded property. \citet{ye2025autonomy} propose four levels, from no autonomy through recommendation and action-with-confirmation to full autonomy, and evaluate models at each. This is a genuine advance, because it makes the object of evaluation explicit. But a level describes capability rather than permission: it says what a system is able to do unsupervised, not what it should be allowed to do, nor who answers when it does it. \citet{habli2020accountability} confront the consequence directly, arguing that accountability in AI-supported care cannot be reduced to the performance of a model and must be designed into the sociotechnical system around it. \citet{clayton2025clinicians} defends the clinician's diagnostic role as something that cannot be outsourced without loss, and \citet{abdulnour2025education} turns the point into a curriculum, proposing that supervision of AI use be taught as a clinical skill in its own right.

The gap here is structural. The levels literature lacks an account of authorization; the accountability literature lacks an operational form. Neither tells a clinician what to write down before allowing an agent to act.

\paragraph{Stream 3: collaboration and its costs.}
A third stream studies the human side, empirically and conceptually. \citet{ong2026flight} import aviation's flight rules into clinical AI and argue that collaboration should be governed by explicit, rehearsed procedures rather than by goodwill. \citet{busch2025attitudes} survey hospital patients across several countries and find that acceptance depends on the task, the degree of human involvement and the clinical stakes, rather than on AI in the abstract. The costs are documented too. In a multicentre observational study, endoscopists exposed to AI-assisted colonoscopy showed an absolute decline of 6.0 percentage points in adenoma detection once the assistance was withdrawn \citep{budzyn2025deskilling}; \citet{natali2025deskilling} generalize the finding into a taxonomy of AI-induced deskilling; and \citet{hunter2024reliance} propose reliance drills as a way of measuring dependence before it becomes dangerous. The lineage runs back to \citet{bainbridge1983}, whose ironies of automation --- that automating a task makes the human's residual role harder and less practised --- describe the present situation with uncomfortable precision. \citet{mollick2024} reframes the relationship as co-intelligence, in which the human sets direction and the machine supplies capacity.

What this stream supplies is a vocabulary of failure. What it does not supply is a rule: it can say when collaboration goes wrong, but not what a legitimate delegation consists of.

\paragraph{Research gap.}
Taken together, the three streams leave a specific hole. One measures capability, one grades autonomy, one catalogues the costs of collaboration. None specifies what a delegation must contain. If a clinician is to authorize an agent to act, four things are at stake and none is currently a design requirement. The scope of the authorization must be stated in advance. The authorization must be revocable while the action is in flight. The action must be reconstructable afterwards. And the human's unaided ability must survive the arrangement. The literature can describe each of these failures separately. It has no account of them as a single object --- and that is the object this paper constructs.

%% =====================================================================
%%  3 / 7   Methodology  方法                                    成员 3
%%  两阶段混合方法：目的性综述 + 情景问卷实验
%% =====================================================================
\section{Methodology}
\label{sec:method}

\paragraph{Design.}
This is a mixed-methods study in two phases. Phase~1 is a purposive review that builds the corpus and derives the themes reported in Section~\ref{sec:results}. Phase~2 is a vignette experiment, specified in Section~\ref{sec:protocol}, that tests the advisor-to-agent shift on human judgements. Phase~1 is described first because the themes it produces determine what Phase~2 has to measure.

\subsection{Phase 1: Purposive Review}
\label{sec:review}

\paragraph{Search strategy.}
We searched PubMed, Web of Science, Scopus, IEEE Xplore, the ACM Digital Library and arXiv for work published between January 2019 and September 2026. The query combined an agent term, a clinical term and a governance term:
\begin{quote}\small
\texttt{("medical agent" OR "clinical agent" OR "agentic AI")\\
AND ("human-in-the-loop" OR oversight OR accountability\\
OR deskilling OR trust)}
\end{quote}
Backward and forward snowballing from the included studies supplemented the database search. We do not claim a systematic review and report no PRISMA diagram: the aim was a defensible and reproducible corpus, not an exhaustive one, and the coding below is reported so that a reader can contest it.

\paragraph{Selection criteria.}
Studies were included if they (i) concern a system that plans, uses tools or executes multi-step clinical tasks, (ii) discuss the human role, oversight, accountability or trust, and (iii) are peer-reviewed or a preprint with a documented method. They were excluded if they (i) evaluate single-turn question answering with no action, (ii) report no analysis of the human--AI arrangement, or (iii) lie outside clinical medicine. Seventeen studies were retained; they are listed in Table~\ref{tab:corpus}.

\paragraph{Extraction and coding.}
Each study was coded on four fields: the core claim; the locus of action, that is, who acts on the patient or the record; the level of autonomy, using the four-level scheme of \citet{ye2025autonomy}; and the type of evidence, classified as benchmark, clinical, attitude survey or conceptual. Coding was deductive for the autonomy field and inductive for the thematic analysis. A theme was retained only if at least two independent sources supported it.

\subsection{Phase 2: Vignette Experiment}
\label{sec:protocol}

\paragraph{Design.}
A $2\times2$ between-subjects design crosses the role of the system (advisor versus agent) with disclosure (the system is identified as an AI agent versus left unspecified). Participants see four fictional vignettes at one randomly assigned condition, which keeps the vignettes within-subject while preventing contamination across conditions.

\paragraph{Materials.}
Four vignettes were written to span clinical risk: oral hypoglycaemic adjustment (low), withholding an ACE inhibitor for hyperkalaemia with an urgent electrocardiogram (moderate), automatic triage escalation for chest pain (moderate to high), and an automated adverse-drug-reaction notice sent to a patient (high). Each vignette exists in four versions that differ only in the final sentence, which states either that the system recommends an action or that it has already carried it out, and either names the system as an AI agent or leaves the actor unspecified.

\paragraph{Measures.}
A manipulation check asks who completed the action. Trust is measured with three seven-point items, willingness to delegate with two (one reverse-scored), responsibility attribution with a single forced-choice item, and perceived risk with one item. Background variables are collected last.

\paragraph{Analysis plan.}
Trust items are averaged after internal consistency is confirmed (Cronbach's $\alpha>0.7$). Effects are tested with a two-way analysis of variance on role and disclosure, and responsibility attribution with a chi-square test of independence. If fewer than forty valid responses are obtained, the design falls back to independent-samples $t$ tests per vignette, reported with Cohen's $d$.

\paragraph{Ethics.}
Participation is voluntary and anonymous. No real patient data are used and all vignettes are fictional. The protocol is specified in full here so that the design can be assessed independently of what it eventually measures.

%% =====================================================================
%%  4 / 7   Results  结果                                        成员 4
%%  主要任务：呈现真实得到的结果
%%    a. 只放测到的数字（准确率/分数/耗时…），解释留给第 5 部分
%%    b. 主结果：一张表 + 一张图，正文里每张表/图都要被 \ref 引用一次
%%    c. 消融 / 敏感性 / 个案分析各放一个 \subsection
%%  篇幅建议：2--3 页；表格用 booktabs + \small，列间距 \tabcolsep 4.8--8pt
%% =====================================================================
\section{Results}
\label{sec:results}

\todo{PART 4 --- Results (Member 4): report what was actually measured, in tables and figures, without interpretation.}

% ---- 主结果表模板（换成你们的数据）----
%
% \begin{table}[t]
% \centering\small
% \setlength{\tabcolsep}{6pt}
% \begin{tabular}{@{}lccc@{}}
% \toprule
% Method & Metric A & Metric B & Metric C \\
% \midrule
% Baseline            & 00.0 & 00.0 & 00.0 \\
% \rowcolor{oursrow}[\tabcolsep][\tabcolsep]Ours & 00.0 & 00.0 & 00.0 \\
% \bottomrule
% \end{tabular}
% \caption{\textbf{One-sentence takeaway of the table.} Setting, metric
% definition, and how many runs are averaged. Best value per column in bold.}
% \label{tab:main}
% \end{table}

% ---- 建议小节 ----
%
% \subsection{Main Results}
% \label{sec:main}
% 先给结论句，再指向 Table~\ref{tab:main} 和对应的图。
%
% \subsection{Ablation Study}
% \label{sec:ablations}
% 每次只改一个因素，逐行累加时在表格里用 "+" 标记。
%
% \subsection{Case Study}
% \label{sec:cases}
% 一个成功、一个失败，各引一小段原文/日志。

%% =====================================================================
%%  5 / 7   Discussion  讨论                                    成员 5
%%  主要任务：解释结果、对照、意义和局限
%%    a. 机制：为什么会出现第 4 部分的结果（不只是复述数字）
%%    b. 对照：与第 2 部分梳理的已有工作一致/不一致在哪，为什么
%%    c. 意义：理论意义与实践意义各一句
%%    d. 局限：数据、规模、评价方式、假设；说清这些局限对结论的威胁
%%  篇幅建议：1.5--2 页；不要引入第 4 部分没有的新结果
%% =====================================================================
\section{Discussion}
\label{sec:discussion}

\todo{PART 5 --- Discussion (Member 5): explain the results, compare them with prior work, and state implications and limitations.}

% ---- 写作骨架 ----
%
% \paragraph{Why the results look like this.}
% 把第 4 部分的现象归因到 2--3 个机制上，每个机制给一句证据（可回引图表）。
%
% \paragraph{Relation to prior work.}
% 与 \ref{sec:litreview} 中的主线逐条对照：哪里一致、哪里相反、差异来源。
%
% \paragraph{Implications.}
% 对后续研究/工程实践各一句话，不要空泛。
%
% \paragraph{Limitations.}
% 明确列出 2--4 条，并说明它们如何限制结论的适用范围。

%% =====================================================================
%%  6 / 7   Conclusion  结论                                    成员 6
%% =====================================================================
\section{Conclusion}
\label{sec:conclusion}

This paper began from a shift in what clinical AI does: systems that once returned a judgement now execute actions, while the evaluation literature continues to grade them on the judgement. We asked what changes when the unit of analysis moves from the advice to the action. A review of seventeen studies across engineering and human-factors work produced four themes --- deskilling, the supervision illusion, accountability diffusion, and trust under uncertainty --- and the central result is that these are not four problems. They are four views of one missing object: the authorization boundary that would state what an agent may do, allow that permission to be withdrawn, leave a record behind, and preserve a human able to act without it.

The answer to the question posed in Section~\ref{sec:intro} is therefore that the human--AI relationship in medicine is not primarily a matter of trust, and cannot be settled by improving accuracy. It is a matter of authority: who may act, within what scope, on whose permission, and with what residue of responsibility when the action goes wrong. Governed delegation makes those questions answerable in advance rather than after an incident, and per delegation rather than per system.

Three directions follow. The vignette experiment in Section~\ref{sec:protocol} should be run, so that the advisor-to-agent shift is measured rather than assumed; the four requirements should be operationalised into an audit schema a deployment review can check; and the retention requirement should be tested against real training programmes, since its violation is invisible until the assistance is removed. What the field needs next is not a better benchmark but a better contract.


\ifshowbib
  \clearpage
  \bibliographystyle{unsrtnat}
  \bibliography{main}
\fi

%% ---- 需要附录时解除注释（证明、完整协议、额外表格）----
% \clearpage
% \appendix
% \input{sections/07-appendix}

%% =====================================================================
%%  成员 7：统稿检查清单（交稿前逐条过一遍）
%%   [ ] 七个部分齐全，章节编号连续，没有人漏交
%%   [ ] 编译日志里没有 "??"（\ref / \citep 全部解析）
%%   [ ] 图表标题格式统一：\textbf{结论句.} + 设置 + 指标口径
%%   [ ] 术语、缩写、\method 写法全文一致；摘要里也统一
%%   [ ] 正文 \citep 与 main.bib 一一对应，没有缺项也没有多余条目
%%   [ ] 参考文献按出现顺序（unsrtnat），作者名/年份/会议名核对过
%%   [ ] \todo 全部清空；页边距、字号未被局部改动
%% =====================================================================

\end{document}
